\documentclass[11pt]{article}

\usepackage[margin=1in]{geometry}
\usepackage[T1]{fontenc}
\usepackage[utf8]{inputenc}
\usepackage{amsmath,amssymb}
\usepackage{booktabs}
\usepackage{tabularx}
\usepackage{array}
\usepackage{enumitem}
\usepackage{xurl}
\usepackage[hidelinks]{hyperref}
\emergencystretch=3em

\title{Supplementary Methods and Provenance\\
A Pre-Individuation Structural Comparison Procedure for Cognitive-Capacity Claims}
\author{Rintaro Sakamoto}
\date{}

\begin{document}
\maketitle
\sloppy

\section{Purpose and evidential boundary}

This supplement records pre-specification chronology, immutable repository authority, source-adjudication boundaries, comparator definitions, and deterministic audit rules that are intentionally compressed in the main article. These materials establish inspectable provenance and internal consistency. They do not constitute independent empirical validation, an independent second coding, or a claim that the repository state is itself philosophical evidence.

The direct analytical objects are source-grounded scientific claims about cognitive capacities. The scientific invariants carried into the present revision are:
\begin{itemize}[leftmargin=*]
\item whole-claim comparison: 1,770/1,770 incomparable;
\item bounded reconstruction: 206 reusable structural subobjects and 99 cross-stratum families;
\item reverse projection: 21 full, 22 partial, and 17 residual;
\item residuals requiring a new top-level role: 0/17;
\item prospective 40-paper result: A0=40, A1=0, A2=0;
\item component arm: 24/24 A0; integrated arm: 16/16 A0;
\item global genealogical independence: not claimed;
\item population or prevalence inference: not licensed;
\item independent human re-adjudication: not performed.
\end{itemize}

\section{Revision authority and freeze chronology}

The present revision starts from exact manuscript/research HEAD
\[
\texttt{095c6f41bc7ca6a08bf423137bbf3b8f7bc068e5}.
\]
That state is the final M10 authority and is the exact parent of the M11 branch.

\begin{table}[htbp]
\centering
\small
\caption{M11 freeze and execution chronology.}
\begin{tabularx}{\textwidth}{>{\raggedright\arraybackslash}p{0.19\textwidth} >{\raggedright\arraybackslash}p{0.29\textwidth} >{\raggedright\arraybackslash}X}
\toprule
Commit & Function & Boundary \\
\midrule
\path{a212820d225050bbc10d395685e74cd2fb0d9912} &
Freeze prospective comparator and CPCG specifications &
No comparator-result or CPCG-result artifact existed in this commit. \\
\path{b31f670f31d2f015b65d8bbf2e4fb12178b46694} &
Pre-result CPCG clarification &
Maps auxiliary \texttt{axis=null} nodes to generic CPCG family \(N\); no result had yet been materialized. \\
\path{1c22fcb177c434d2900ffb3af9ef6baee2af73ff} &
Execute frozen direct-preservation and CPCG controls &
Materializes direct-preservation comparator, CPCG, and same-label expository artifacts. \\
\path{5769537203b452e0c86c69fabf9a7d3240ca4909} &
Freeze encoding-permissive composition baseline &
Separates the original A-state composition semantics from direct distinction preservation; no result artifact yet exists. \\
\path{38d4b5a13e04127a945d31374924dd089cfffca6} &
Execute encoding-permissive baseline &
Materializes the literal weak-basis A-state comparison under unchanged source mechanisms. \\
\bottomrule
\end{tabularx}
\end{table}

The comparator freeze explicitly records two non-blind facts known before result materialization. First, the frozen prospective matrix already showed all eight working-basis coordinates in all 40 accepted rows, so some closure outcomes were logically anticipatable. Second, candidate same-label pairs had been explored before the freeze; the MEM04--CNC05 example is therefore expository rather than a predeclared inferential test.

\section{Original 60-claim staging}

The 60-claim analysis used reviewed source-grounded structured records. Historical construct labels, author identity, institution, venue, and citation count were excluded from structural similarity logic. The working basis, whole-claim projection, and bounded-reconstruction operator were fixed before later role-grammar interpretation.

The universal whole-claim failure was preserved rather than rewritten:
\[
\binom{60}{2}=1770,\qquad 1770/1770\ \text{incomparable}.
\]
The subsequent bounded reconstruction changed the unit of reuse rather than the earlier verdict. It recovered 206 bounded objects, of which 99 had cross-stratum support. The later reverse projection was performed on the fixed records and produced 21 full, 22 partial, and 17 residual cases. Residual adjudication found no case requiring an additional top-level cognitive-system role.

Machine-readable authority includes:
\begin{itemize}[leftmargin=*]
\item \path{research/paper2/reference_global_phi_freeze_v1.json};
\item \path{research/paper2/bounded_xlike_reconstruction_v1.json};
\item \path{research/paper2/grammar_v0_reverse_projection_aggregate_v1.json}.
\end{itemize}

\section{Prospective 40-paper source and adjudication workflow}

The prospective denominator contains exactly 40 cognitive-model papers: 24 component-oriented papers and 16 integrated papers. No unfavorable result could trigger source substitution. Exact DOI overlap with the original 60-claim corpus is zero, but DOI separation is not treated as genealogy independence.

Each paper retained source-first decomposition and a paper-level architectural outcome. The aggregate frozen matrix is
\path{research/paper2/p399/main/integration/M7/M7_40_PAPER_ARCHITECTURAL_MATRIX_v1.json}.
The paper-level evidence table is
\path{research/paper2/p399/main/integration/M9B/M9B_MAIN40_PAPER_LEVEL_EVIDENCE_v1.json}.

The least-reconstruction decision variable is:
\[
A0=E_0,\qquad
A1=E_1\setminus E_0,\qquad
A2=E_2\setminus E_1,
\]
where \(E_0\subseteq E_1\subseteq E_2\) are source contracts faithfully realizable with, respectively, source-defined mechanisms only, reconstruction-added stateless mediation, and reconstruction-added persistent/stateful mediation.

Source-defined hierarchy, recurrence, gating, planner state, memory, or belief state is compatible with A0. A1 and A2 count only reconstruction-added structure. The certified prospective result remains immutable:
\[
A0=40,\qquad A1=0,\qquad A2=0.
\]

\section{Two prospective representation baselines}

The central hostile-review question has two readings. The first uses the original A0/A1/A2 decision variable: does a weaker representation require a \emph{new mechanism} once source-defined mechanisms remain available? The second uses the stricter direct-preservation objective: does a weaker representation keep the selected distinctions first-class without explicit re-exposure? These are not equivalent.

\subsection{Encoding-permissive composition baseline}

The tested arms are the full working basis, the dynamic backbone \(\{X,K,T\}\), the POMDP-like view \(\{X,P,K,Q,T,\rho/O\}\), and four single-role ablations removing \(\Pi\), \(C\), \(P\), or \(T\).

Under the original A-state semantics, representation-only recoding is not a reconstruction-added mechanism. The frozen encoding contract permits source-defined boundaries, constraints, criteria, interfaces, observations, and temporal markers to remain scientifically unchanged while being represented inside more generic state/transition/interface structure. For example, a boundary can be carried as source-defined typing on \(X\), a criterion as source-defined parameterization of \(X/K\), and an observation as a typed carrier plus source-defined readout mapping. Such recoding may destroy first-class role identity, but it does not add a coordinator, adapter, or persistent state.

The result is:
\begin{table}[htbp]
\centering
\caption{Encoding-permissive A-state closure.}
\begin{tabular}{lrrr}
\toprule
Arm & A0 & A1 & A2 \\
\midrule
Full & 40 & 0 & 0 \\
Dynamic & 40 & 0 & 0 \\
POMDP-like & 40 & 0 & 0 \\
Minus \(\Pi\) & 40 & 0 & 0 \\
Minus \(C\) & 40 & 0 & 0 \\
Minus \(P\) & 40 & 0 & 0 \\
Minus \(T\) & 40 & 0 & 0 \\
\bottomrule
\end{tabular}
\end{table}

This confirms the reviewer-level triviality concern in the relevant sense: prospective 40/40 A0 does not discriminate the exact comparison basis. It supports the narrower claim that no reconstruction-added coordination mechanism is required for the 40 source models under any tested encoding-capable view.

Machine-readable authority:
\begin{itemize}[leftmargin=*]
\item \path{research/paper2/p399/main/integration/M11/M11_ENCODING_PERMISSIVE_COMPOSITION_BASELINE_SPEC_v1.json};
\item \path{research/paper2/p399/main/integration/M11/M11_ENCODING_PERMISSIVE_COMPOSITION_BASELINE_RESULTS_v1.json}.
\end{itemize}

\subsection{Direct-preservation closure control}

The separately frozen direct-preservation control uses the same arms but charges preservation debt when a declared source distinction is no longer first-class. A0 requires no preservation debt. A1 is assigned when every omitted distinction can be re-exposed by a load-bearing but stateless structure whose semantic content is already source-defined. This contract forbids hiding the distinction inside an opaque generic state/transition encoding.

\begin{table}[htbp]
\centering
\caption{Direct-preservation closure result on the same frozen 40 rows.}
\begin{tabular}{lrrr}
\toprule
Arm & A0 & A1 & A2 \\
\midrule
Full & 40 & 0 & 0 \\
Dynamic & 0 & 40 & 0 \\
POMDP-like & 0 & 40 & 0 \\
Minus \(\Pi\) & 0 & 40 & 0 \\
Minus \(C\) & 0 & 40 & 0 \\
Minus \(P\) & 0 & 40 & 0 \\
Minus \(T\) & 0 & 40 & 0 \\
\bottomrule
\end{tabular}
\end{table}

The two baselines jointly delimit the inference. Generic encodability is broad enough that all tested views can remain A0, but only the full working basis keeps the declared distinctions directly available without stateless re-exposure. The prospective no-added-coordinator result and the basis-selection argument are therefore supported by different evidence.

Direct-preservation machine-readable authority:
\begin{itemize}[leftmargin=*]
\item \path{research/paper2/p399/main/integration/M11/M11_PROSPECTIVE_COMPARATOR_SPEC_v1.json};
\item \path{research/paper2/p399/main/integration/M11/M11_PROSPECTIVE_COMPARATOR_RESULTS_v1.json}.
\end{itemize}

\section{Representation-sensitivity controls}

\subsection{TAIF: lossless non-uniqueness}

The Tagged Admissibility--Interface Factorization maps the original role distinctions into five first-class families \(\{X,A,I,K,T\}\), retaining subtype tags that make the map invertible. It round-trips the nine role atoms, all eight working-basis axes, all 60 reverse projections, and all 206 bounded objects. Its scientific role is to demonstrate that the exact first-class factorization is not unique. Because it is invertible, it is not treated as the principal robustness test.

\subsection{CPCG: non-invertible coarse process representation}

The Coarse Process--Constraint Graph uses:
\[
\{N,G,I,E,T\},
\]
where \(N\) merges state/carrier and partition/boundary roles, \(G\) merges constraints and criteria, \(I\) merges probe/intervention and observation/readout, \(E\) collapses non-temporal transformation/dependence kinds, and \(T\) collapses detailed temporal subtypes to temporal presence/order.

The projection deliberately discards fine role strings, control subtype names, non-temporal relation kinds, and temporal subtype distinctions. It is therefore non-injective.

The frozen result is:
\begin{itemize}[leftmargin=*]
\item 206 bounded objects \(\rightarrow\) 143 CPCG signatures;
\item 37 CPCG collision groups among originally distinct archetypes;
\item 63/99 original cross-stratum families collide with another family;
\item 62 CPCG signatures still have cross-stratum support;
\item 58/60 claims participate in at least one cross-stratum CPCG signature;
\item the ATT01--BLF01 different-label/shared-structure witness remains shared.
\end{itemize}

The result is descriptive sensitivity under one lossy rival. It is not representation neutrality and does not identify a true five-family ontology.

Machine-readable specification and result:
\begin{itemize}[leftmargin=*]
\item \path{research/paper2/p399/main/integration/M11/M11_LOSSY_RIVAL_SPEC_v1.json};
\item \path{research/paper2/p399/main/integration/M11/M11_LOSSY_RIVAL_PREFREEZE_AMENDMENT_v1.json};
\item \path{research/paper2/p399/main/integration/M11/M11_LOSSY_RIVAL_RESULTS_v1.json}.
\end{itemize}

\section{Two-way label witness}

The different-label direction is illustrated by ATT01 and BLF01, which retain a shared bounded object despite different historical labels. The converse direction is illustrated by MEM04 and CNC05. Both frozen records include the historical label \emph{semantic memory}, but their bounded support sets are disjoint. MEM04 supports partition/state/temporal and history-indexed state objects; CNC05 supports state-to-outcome dependence objects. Their CPCG signature sets also remain disjoint.

This pair was selected after same-label candidates had been inspected and is therefore an expository witness, not a predeclared prevalence or hypothesis test. Authority:
\path{research/paper2/p399/main/integration/M11/M11_SAME_LABEL_WORKED_EXAMPLE_v1.json}.

\section{Single-adjudicator and reproducibility boundary}

Independent human re-adjudication was not performed. Language-model checks, Lean proofs, deterministic continuous integration, synthetic reachability tests, real-source counterfactual perturbations, and comparator controls test different properties. None is counted as an independent human coder.

The supported claim is:
\begin{quote}
Procedural auditability is established; inter-rater reliability remains unmeasured.
\end{quote}

Accordingly, exact counts should be read as results under the frozen source review and structural adjudication, not as estimates of inter-rater agreement. External re-adjudication remains a separate validation opportunity.

\section{Repository and continuous-integration role}

Deterministic scripts verify arithmetic, immutable scientific invariants, comparator counts, CPCG counts, manuscript claim boundaries, and strict LaTeX compilation. Continuous integration establishes that the reported artifacts are mutually consistent at an exact repository state. It does not establish empirical adequacy or philosophical correctness.

The final revision workflow is
\path{.github/workflows/p399-main-m11-discriminative-baselines.yml}.
Its audit script is
\path{scripts/paper2_m11_discriminative_baselines_audit.py}.

\end{document}
